% SEGMENT OLP-0617-S01
% Part: methods
% Chapter: induction
% Section: inductive-definitions

\documentclass[../../../include/open-logic-section]{subfiles}

\begin{document}

\olfileid{mth}{ind}{idf}

\olsection{தொகுத்தறிதல் வரையறைகள்}

தருக்கவியலில் சில வகைப் பொருள்களை நாம் அடிக்கடி
\emph{தொகுத்தறிதல் முறையில்} வரையறுக்கிறோம்.
அதாவது, ஏற்கெனவே அந்த வகையைச் சேர்ந்த
பொருள்களிலிருந்து புதிய பொருள்களை எப்படிப் பெறுவது
என்பதை விளக்கும் விதிகளை, அந்த வகையில் எந்தப்
பொருள் சேரும் என்பதைத் தீர்மானிக்கக் குறிப்பிடுகிறோம்.
எடுத்துக்காட்டாக, ஒரு மொழியின் சொற்களும்
!!{formula}s உம் போன்ற தனிப்பட்ட குறித்தொடர்களை
இவ்வாறு வரையறுக்கிறோம். எளிய எடுத்துக்காட்டாக,
$\mathrm{a}$, $\mathrm{b}$, $\mathrm{c}$,
$\mathrm{d}$ என்ற எழுத்துகள்,~$\circ$ என்ற
குறி, $[$ மற்றும் $]$ என்ற அடைப்புகள்
ஆகியவற்றைக் கொண்ட ``$[[\mathrm{c} \circ \mathrm{d}][$'',
``$[\mathrm{a}[]\circ]$'', ``$\mathrm{a}$''
அல்லது ``$[[\mathrm{a} \circ \mathrm{b}]\circ
\mathrm{d}]$'' போன்ற தொடர்களைப் பாருங்கள்.
முதல் இரண்டு தொடர்களிலும் ஏதோ ``தவறு'' இருப்பதாக
உங்களுக்குத் தோன்றலாம்: முதலாவதில் அடைப்புகள்
சற்றும் ``பொருந்தவில்லை'';~``$\circ$'' என்பதும்
தனித்தனியே பொருள் தரும் கோவைகளை ``இணைக்க''
வேண்டும் என்று நினைக்கலாம். மூன்றாவதும் நான்காவதும்
மேலானவையாகத் தோன்றுகின்றன: ஒவ்வொரு ``$[$''
க்கும் உரிய மூடும் ``$]$'' உள்ளது (திறக்கும்
அடைப்பு ஏதேனும் இருந்தால்); ஒவ்வொரு $\circ$
வின் இருபுறமும் ஒரு ஜோடி அடைப்புக்குள்
``ஒழுங்கான'' கோவைகள் உள்ளன.

% SEGMENT OLP-0617-S02
``ஒழுங்கான சொல்'' என்றால் துல்லியமாக எது
என்று குறிப்பிட விரும்புகிறோம். முதலில்,
தனியே உள்ள ஒவ்வோர் எழுத்தும் ஒழுங்கானது.
ஒரு தனி எழுத்தாக இல்லாத எதுவும்
``$[t \circ s]$'' என்ற வடிவில் இருக்க வேண்டும்;
இங்கே $s$, $t$ இரண்டும் தாமே ஒழுங்கானவை.
மறுதிசையில், $t$, $s$ ஒழுங்கானவை எனில்,
அவற்றுக்கு இடையில் $\circ$ ஐ வைத்து
இருபுறமும் அடைப்புகளைச் சேர்த்து புதிய
ஒழுங்கான சொல்லை உருவாக்கலாம். இச்செயல்களைப்
பயன்படுத்தி ஒழுங்கான சொற்களின் கணத்தை
\emph{வரையறுக்கலாம்}. இது ஒரு
\emph{தொகுத்தறிதல்
  வரையறை}.

% SEGMENT OLP-0617-S03
\begin{defn}[ஒழுங்கான சொற்கள்]
  \emph{ஒழுங்கான சொற்களின்} கணம் பின்வருமாறு
  தொகுத்தறிதல் முறையில் வரையறுக்கப்படுகிறது:
  \begin{enumerate}
  \item $\mathrm{a}$, $\mathrm{b}$, $\mathrm{c}$,
    $\mathrm{d}$ ஆகிய எழுத்துகள் ஒவ்வொன்றும்
    ஒழுங்கான சொல்.
  \item $s_1$, $s_2$ ஒழுங்கான சொற்களெனில்,
    $[s_1 \circ s_2]$ உம் ஒழுங்கான சொல்.
  \item வேறெதுவும் ஒழுங்கான சொல் அல்ல.
  \end{enumerate}
\end{defn}

% SEGMENT OLP-0617-S04
இந்த வரையறை, (1),~(2) என்ற இரு நிபந்தனைகளின்படி
முடிவுறு எண்ணிக்கையிலான படிகளில் உருவாக்க முடிந்தால்,
முடிந்தால் மட்டுமே, ஒன்று ஒழுங்கான சொல் என்று
கூறுகிறது. முதல் படியில் தனித்தனி எழுத்துகளான
ஒழுங்கான சொற்கள் அனைத்தையும் உருவாக்குகிறோம்:
\[
\mathrm{a}, \mathrm{b}, \mathrm{c}, \mathrm{d}
\]
இரண்டாம் படியில், உருவாக்கிய சொற்களுக்கு (2) ஐப்
பயன்படுத்துகிறோம். இரு எழுத்துகளை இணைக்கும்
எல்லா வழிகளுக்கும் பின்வருவன கிடைக்கும்:
\[
[\mathrm{a} \circ \mathrm{a}], [\mathrm{a} \circ \mathrm{b}],
[\mathrm{b} \circ \mathrm{a}], \dots, [\mathrm{d} \circ \mathrm{d}]
\]
மூன்றாம் படியில், இதுவரை உருவாக்கிய ஏதேனும்
இரண்டு ஒழுங்கான சொற்களுக்கு (2) ஐ மீண்டும்
பயன்படுத்துகிறோம். எடுத்துக்காட்டாக,
$[\mathrm{a} \circ [\mathrm{a} \circ
    \mathrm{a}]]$ கிடைக்கும்; இதில் $t$ என்பது
முதல் படியில் கிடைத்த $\mathrm{a}$, $s$
என்பது இரண்டாம் படியில் கிடைத்த
$[\mathrm{a} \circ \mathrm{a}]$. அதேபோல்,
இரண்டாம் படியின் $[\mathrm{b} \circ \mathrm{c}]$,
$[\mathrm{d} \circ \mathrm{b}]$ என்ற இரு
சொற்களிலிருந்து $[[\mathrm{b} \circ
    \mathrm{c}] \circ [\mathrm{d} \circ \mathrm{b}]]$
கிடைக்கும். இப்படியே தொடரலாம். இவ்வாறு
உருவாக்கப்படாத எதுவும் ஒழுங்கான சொற்களின்
கணத்திற்குள் வந்துவிடாதபடி (3) தடுக்கிறது.

% SEGMENT OLP-0617-S05
உணர்வுக்கு ``ஒழுங்காகத்'' தோன்றும் ஒவ்வொரு
குறித்தொடரும் இந்த வரையறைப்படி ஒழுங்கானது
என்று நாம் இன்னும் நிறுவவில்லை என்பதைக்
கவனியுங்கள். ஆனால் நாம் உருவாக்கக்கூடிய
ஒவ்வொன்றும் உண்மையிலேயே ``ஒழுங்காகத்''
தோன்றுவது தெளிவு: அடைப்புகள் பொருந்துகின்றன;
$\circ$ தாமே ஒழுங்கான பகுதிகளை இணைக்கிறது.

% SEGMENT OLP-0617-S06
தொகுத்தறிதல் வரையறைகளின் முக்கிய அம்சம்:
எல்லா ஒழுங்கான சொற்களையும் பற்றிய கூற்றை
நிறுவ வேண்டுமானால், எந்த நிகழ்வுகளை ஆராய
வேண்டும் என்று வரையறையே சொல்கிறது.
எடுத்துக்காட்டாக, $t$ ஓர் ஒழுங்கான சொல்
என்று தெரிந்தால், $t$ இன் வடிவம் என்னவாக
இருக்கலாம் என்று வரையறை காட்டுகிறது:
$t$ ஓர் எழுத்தாக இருக்கலாம்; அல்லது
ஒழுங்கான சொற்கள் $s_1$,~$s_2$ எவற்றுக்கோ
$[s_1 \circ s_2]$ ஆக இருக்கலாம்.
(3) என்ற கூறினால் இவ்விரண்டைத் தவிர
வேறு வாய்ப்பில்லை.

% SEGMENT OLP-0617-S07
தொகுத்தறிதல் முறையில் வரையறுக்கப்பட்ட கணத்தின்
எல்லா உறுப்புகளையும் பற்றிய கூற்றுகளை
நிறுவும்போது, தொகுத்தறிதலின் வலுவான வடிவம்
மிகவும் முக்கியமாகிறது. எடுத்துக்காட்டாக,
நீளம்~$n$ உடைய ஒவ்வோர் ஒழுங்கான சொல்லிலும்
$[$ களின் எண்ணிக்கை~$< n/2$ என்பதை
நிறுவ விரும்புகிறோம் என்க. எல்லா~$n$
களையும் பற்றிய கூற்றாக இதைக் கொள்ளலாம்:
ஒவ்வோர் $n$ க்கும், நீளம்~$n$ உடைய
எந்த ஒழுங்கான சொல்லிலும் $[$ களின்
எண்ணிக்கை $< n/2$.

% SEGMENT OLP-0617-S08
\begin{prop}
  எந்த $n$ க்கும், நீளம்~$n$ உடைய
  ஓர் ஒழுங்கான சொல்லில் $[$ களின்
  எண்ணிக்கை $< n/2$.
\end{prop}

% SEGMENT OLP-0617-S09
\begin{proof}
இந்த முடிவை (வலுவான) தொகுத்தறிதலால் நிறுவ,
பின்வரும் நிபந்தனைக் கூற்று மெய் என்று காட்ட வேண்டும்:
\begin{quote}
  எல்லா $l < k$ க்கும், நீளம்~$l$ உடைய
  ஒவ்வோர் ஒழுங்கான சொல்லிலும் $[$ களின்
  எண்ணிக்கை $< l/2$ எனில், நீளம்~$k$
  உடைய ஒவ்வோர் ஒழுங்கான சொல்லிலும் $[$ களின்
  எண்ணிக்கை $< k/2$.
\end{quote}
இந்த நிபந்தனையைக் காட்ட, அதன் முற்பகுதி மெய்
என்று எடுத்துக்கொள்க: எந்த $l<k$ க்கும்,
நீளம்~$l$ உடைய ஒழுங்கான சொற்களில்
$[$ களின் எண்ணிக்கை $< l/2$ என்க.
இதையே தொகுத்தறிதல் எடுகோள் என்கிறோம்.
நீளம்~$k$ உடைய ஒழுங்கான சொற்களுக்கும்
அதே முடிவைக் காட்ட விரும்புகிறோம்.

% SEGMENT OLP-0617-S10
$t$ என்பது நீளம்~$k$ உடைய ஒழுங்கான
சொல் என்க. ஒழுங்கான சொற்கள் தொகுத்தறிதல்
முறையில் வரையறுக்கப்பட்டதால் இரு நிகழ்வுகள்:
(1)~$t$ தனி எழுத்து; அல்லது (2)~$t$
என்பது சில ஒழுங்கான சொற்கள் $s_1$,~$s_2$
க்கு $[s_1 \circ s_2]$.
\begin{enumerate}
\item $t$ ஓர் எழுத்து. அப்போது $k = 1$;
  $t$ இல் உள்ள $[$ களின் எண்ணிக்கை~$0$.
  $0 < 1/2$ என்பதால் கூற்று மெய்.
\item $t$ என்பது சில ஒழுங்கான சொற்கள்
  $s_1$,~$s_2$ க்கு $[s_1 \circ s_2]$.
  ~ $s_1$ இன் நீளம் $l_1$,~$s_2$ இன்
  நீளம் $l_2$ என்க. அப்போது $t$ இன்
  நீளம்~$k$ என்பது $l_1+l_2+3$; அதாவது
  $s_1$, $s_2$ இன் நீளங்களுடன் $[$,
  $\circ$, $]$ என்ற மூன்று குறிகளையும்
  சேர்த்தது. $l_1+l_2+3$ எப்போதும்
  $l_1$ ஐவிடப் பெரியது; ஆகவே $l_1 < k$.
  அதேபோல் $l_2 < k$. எனவே தொகுத்தறிதல்
  எடுகோள் $s_1$,~$s_2$ இரண்டுக்கும்
  பொருந்தும்: $s_1$ இல் உள்ள $[$ களின்
  எண்ணிக்கை~$m_1$ என்பது $< l_1/2$;
  ~ $s_2$ இல் உள்ள $[$ களின்
  எண்ணிக்கை~$m_2$ என்பது $< l_2/2$.

  $t$ இல் உள்ள $[$ களின் எண்ணிக்கை,
  ~ $s_1$ இல் உள்ள $[$ களின் எண்ணிக்கை,
  ~ $s_2$ இல் உள்ள $[$ களின் எண்ணிக்கை,
  மேலும்~$1$ ஆகியவற்றின் கூட்டுத்தொகை;
  அதாவது $m_1 + m_2 + 1$. $m_1
  < l_1/2$ மற்றும் $m_2 < l_2/2$
  என்பதால்:
  \[
  m_1 + m_2 + 1 < \frac{l_1}{2} + \frac{l_2}{2} + 1 = \frac{l_1+l_2+2}{2} < \frac{l_1+l_2+3}{2} = k/2.
  \]
\end{enumerate}
இரு நிகழ்வுகளிலும், $t$ இல் உள்ள
$[$ களின் எண்ணிக்கை $< k/2$ என்று
(தொகுத்தறிதல் எடுகோளின் அடிப்படையில்)
காட்டிவிட்டோம். வலுவான தொகுத்தறிதலால்
முன்மொழிவு பின்தொடர்கிறது.
\end{proof}

% SEGMENT OLP-0617-S11
\begin{prob}
  மிக ஒழுங்கான சொற்களின் கணத்தைப்
  பின்வருமாறு வரையறுக்கவும்:
  \begin{enumerate}
  \item $\mathrm{a}$, $\mathrm{b}$, $\mathrm{c}$,
    $\mathrm{d}$ ஆகிய எழுத்துகள் ஒவ்வொன்றும்
    மிக ஒழுங்கான சொல்.
  \item $s$ மிக ஒழுங்கான சொல் எனில்,
    $[s]$ உம் மிக ஒழுங்கான சொல்.
  \item $s_1$, $s_2$ மிக ஒழுங்கான சொற்களெனில்,
    $[s_1 \circ s_2]$ உம் மிக ஒழுங்கான சொல்.
  \item வேறெதுவும் மிக ஒழுங்கான சொல் அல்ல.
  \end{enumerate}
  நீளம்~$n$ உடைய மிக ஒழுங்கான சொல்~$t$
  இல் $[$ களின் எண்ணிக்கை $\le n/2 +1$
  என்று காட்டுக.
\end{prob}

\end{document}
